
\section*{Part V: The Portfolio of Possible Businesses}
\addcontentsline{toc}{section}{Part V: The Portfolio of Possible Businesses}

\subsection{27. Think in engines, not a conglomerate of unrelated courses}

There are at least 5 viable engines around a technical institute:

\begin{enumerate}[leftmargin=*,itemsep=1.5pt]
  \item \textbf{Learner tuition:} people pay to become more capable.
  \item \textbf{Institutional services:} colleges pay for programs, labs, instructional design, faculty development, assessments, or career support.
  \item \textbf{Enterprise services:} employers pay for technical capability, applied research, or measurable workflow improvement.
  \item \textbf{Tools and IP:} customers license useful educational, assessment, or technical products created through accumulated expertise.
  \item \textbf{Philanthropy, sponsorship, and research funding:} funders support public-value inquiry, scholarships, equipment, and institutional permanence, often subject to restrictions.
\end{enumerate}

A 6th engine - startup equity, patent licensing, and venture returns - can produce outsize results but should be treated as high-variance optional upside, not as a reliable operating budget.

A significant organization may eventually operate all 6. A small founder cannot responsibly run all 6 from day 1.

\subsection{28. The adjacency matrix}

\begin{table*}[t]
\centering
\small
\caption{The Institutional Adjacency Matrix: Asset Reuse, Sales Motions, and Failure Modes}
\label{tab:ch36_adjacency_matrix}
\begin{tabularx}{\textwidth}{p{2.8cm} p{2.6cm} p{2.6cm} p{2.6cm} p{2.6cm} X}
\toprule
\textbf{ Expansion } & \textbf{ Asset reuse } & \textbf{ Sales motion } & \textbf{ Typical cash pattern } & \textbf{ New institutional asset } & \textbf{ Main trap } \\
\midrule
 2nd D2C course  &  teaching brand, learner relationships, pedagogy  &  direct-to-learner  &  paid before/\allowbreak during delivery, possible refunds  &  broader alumni and content  &  catalog sprawl  \\
 Advanced fellowship  &  successful alumni, mentors, labs  &  selective D2C/\allowbreak B2B sponsorship  &  small high-touch cohorts  &  top talent and deep work  &  founder-intensive prestige theater  \\
 College elective  &  curriculum and assessment  &  academic B2B  &  negotiated installments, sometimes late  &  distributed academic reach  &  becoming cheap outsourced faculty  \\
 College center of excellence  &  labs, mentors, assessments  &  long-cycle institutional sale  &  milestone/\allowbreak capex-linked  &  physical proof, faculty alliances  &  stranded equipment, ambiguous ownership  \\
 Enterprise AI L\&D  &  subject expertise, case studies  &  consultative B2B  &  contract milestones  &  access to important workflow problems  &  custom workshop treadmill  \\
 Applied AI consulting  &  technical team, employer trust  &  project-based B2B  &  milestones / retainers  &  casework, engineering reputation  &  services swallowing the institute  \\
 Educational/\allowbreak assessment SaaS  &  repeated internal tools  &  product-led and B2B  &  subscription/\allowbreak usage  &  reusable IP, recurring distribution  &  building a product nobody needs  \\
 Applied research group  &  engineering talent and data partnerships  &  grants, commissioned research  &  milestone, restricted, uncertain  &  publications, researcher reputation  &  dependence on 1 sponsor  \\
 International cohort  &  brand, online teaching  &  global D2C  &  prepaid but refund/\allowbreak currency complexity  &  global learner network  &  high CAC, low local relevance  \\
 Physical studio/\allowbreak lab  &  high-touch pedagogy, community  &  local/\allowbreak hybrid  &  lease and equipment precede utilization  &  embodied community, hands-on proof  &  vanity premises  \\
 Formal degree programs  &  standards, faculty, endowment, property  &  regulated university admission  &  term/\allowbreak semester fees, long commitments  &  academic legitimacy  &  premature university-scale liabilities  \\
\bottomrule
\end{tabularx}
\end{table*}

\textbf{\textbf{[Decision rule]}} The 1st 3 choices should be justified by concrete demand and available operators. The last 3 require different funding and governance logic; do not reach them simply because online revenue hit a round number.

\subsection{29. The 3-bucket strategy, correctly interpreted}

Bayes has a useful long-run hypothesis: \textbf{a bounded core}, \textbf{scalable partnerships}, and \textbf{rare power-law experiments}.

\textbf{Bucket 1: Core educational identity.} Courses, cohorts, fellowships, faculty-led teaching, and assessments. Its primary purpose is reliable learner value and reputation; its economic function is recurring cash flow. Beware of forcing infinite enrollment into a model where instructor attention is scarce.

\textbf{Bucket 2: Partner and enterprise businesses.} College programs, enterprise learning, labs, research services, and institution-facing software. These may grow faster because customers possess distribution and budgets. Their contractual complexity, working capital, and customer concentration need independent oversight.

\textbf{Bucket 3: Options.} Research IP, deep-technology products, incubation, licensing, and selective equity stakes. Their expected outcomes may be power-law distributed, but most experiments will not generate a material return; they need spending caps and separation from learner obligations.

A rough mature \textbf{attention} allocation could eventually resemble 50-60\% core, 30-40\% scalable partnerships, and 5-10\% speculative options. This is not an advisable day-1 headcount or expenditure allocation. In the 1st year, even 10\% of founder time on a venture fund could be a damaging distraction.

\subsection{30. College partnerships: vendor, collaborator, or institutional substrate?}

There are at least 4 levels:

\textbf{Level 1 - Teaching vendor:} Bayes delivers a bounded course or placement workshop. Quick to test; little control; readily replaced.

\textbf{Level 2 - Curriculum partner:} Bayes jointly designs electives, trains college faculty, provides assessment standards, and measures outcomes. More integration and stronger evidence.

\textbf{Level 3 - Applied lab partner:} Bayes and the college run labs, capstones, mentorship, and live industry projects under documented academic and IP policies. Stronger relationships, more delivery obligation.

\textbf{Level 4 - Long-term academic collaborator:} joint research, faculty exchanges, co-developed programs where permitted, and possibly university-approved pathways. Potentially valuable, but governance, credential claims, accreditation, and rights must be precise.

Move up only after the previous level works and the partner behaves well. A college that fails to pay a small pilot on time is not automatically a good partner for a \textnormal{₹}2-crore lab.

A center of excellence should never be primarily a room with logos, donated computers, and ceremonial photographs. Its operating thesis should state what it produces: faculty capability, verified student projects, research outputs, externally used tools, internships, or publications.

\textbf{Contract essentials:} sponsor, budget owner, student access, curriculum ownership, brand usage, data handling, technical infrastructure ownership, equipment maintenance, measurable deliverables, faculty hours, billing triggers, payment security, cancellation, and exit/\allowbreak transfer rights.

\subsection{31. Enterprise AI training: sell transformation, not class hours}

An enterprise may pay high fees for 3 kinds of value: reducing workforce capability risk, improving defined workflows, and gaining trusted expertise in difficult technical decisions. Mere ``AI awareness sessions'' are easily commoditized.

The initial offer should be narrow. For example, an engineering team's internal AI evaluation discipline: assess current workflows, train a team on a concrete failure taxonomy, implement evaluation on a designated workflow, and show before/\allowbreak after performance with appropriate safeguards. The buyer should understand the deliverable before seeing the training syllabus.

A sensible progression:

\begin{itemize}[leftmargin=*,itemsep=1.5pt]
  \item \textbf{Paid discovery} or a short diagnostic.
  \item \textbf{Fixed-scope transformation pilot} with measurable acceptance criteria.
  \item \textbf{Repeat contract} for other teams or recurring evaluation.
  \item \textbf{Reusable training + tooling} if several customers reveal the same underlying problem.
  \item \textbf{Applied research} only when the challenge contains a real unsolved question and funding supports investigation.
\end{itemize}

Do not promise client ROI that Bayes cannot causally attribute. If performance fees are proposed, carefully define counterfactuals, attribution, changing external conditions, data access, and payment dispute resolution.

\subsection{32. Consulting: a bridge into research or a permanent trap}

Consulting has a favorable early property: customers may pay for expertise before the provider has a polished product. But it has an unfavorable property: customers often purchase precisely the scarce judgment of the founder. The 1st can fund a school; the 2nd can prevent the school from forming.

3 design principles keep a boutique practice aligned:

\begin{enumerate}[leftmargin=*,itemsep=1.5pt]
  \item \textbf{Select a research-relevant class of problems}, such as evaluation, retrieval reliability, model risk, forecasting, or decision systems.
  \item \textbf{Price for delegation and documentation}, not heroics. A profitable project after paying fair market rates to a delivery team is better evidence than 1 profitable only because the founder works nights for free.
  \item \textbf{Negotiate reusable, nonconfidential learning rights} wherever appropriate, without exposing client IP or data. Aggregate lessons; never publish protected material as ``research.''
\end{enumerate}

When repeat engagements converge, consider a product or dedicated lab. Until then, call it consulting, not research.

\subsection{33. Software products: when a curriculum tool becomes a business}

Technology is leverage only if it solves a recurring bottleneck. An internal system for diagnostics, technical assessment, teaching support, or project feedback may be unusually well informed by Bayes's own students. That makes it a possible product hypothesis.

Test it in this order:

\begin{enumerate}[leftmargin=*,itemsep=1.5pt]
  \item Use it internally until instructors reliably prefer it to manual work.
  \item Measure actual savings and educational outcomes, not just features shipped.
  \item Ask a few unrelated schools whether they will pay for the same capability.
  \item Pilot with them under realistic integration and support costs.
  \item Only then productize onboarding, security, reporting, reliability, and pricing.
\end{enumerate}

Most internal tools are not standalone SaaS companies. The same code can have positive internal ROI and negative external market economics. A product should not be spun out because it is technically elegant.

\subsection{34. When to expand geographically}

Geographic expansion has 3 forms:

\begin{itemize}[leftmargin=*,itemsep=1.5pt]
  \item \textbf{Remote distribution:} learners or enterprise clients in new cities or countries; low fixed cost, more regulatory and customer-acquisition uncertainty.
  \item \textbf{Partner presence:} faculty visits, college labs, short residencies, demonstration days; modest fixed cost and borrowed local credibility.
  \item \textbf{Owned presence:} long-term classroom, laboratory, campus, or foreign entity; expensive and management-intensive.
\end{itemize}

The rational sequence is often remote 1st, partner 2nd, owned last. The exception is a program whose core learning outcome genuinely requires recurring physical lab time and cannot be obtained through a reliable partner.

Open a permanent local center only when measured local demand, occupancy, service standards, and a credible local operator justify it. If the institute's end-state is 1 major research campus, do not accidentally create a franchise network whose leases and local P\&Ls make eventual consolidation politically or financially impossible.

\subsection{35. The ``Intel Inside'' analogy has limits}

A college partner might adopt a recognizable Bayes technical standard, assessment framework, and faculty support layer. This could create a B2B2C distribution and legitimacy engine. But the economics of an embedded technology standard differ from those of outsourced teaching.

To approach a genuine platform position, Bayes would need valuable recurring assets: trusted assessments, instructional infrastructure, recognized quality assurance, and data rights appropriately governed. Colleges would have to benefit from continuing to use Bayes rather than merely repurchasing hours of instruction. Students and employers would need to understand and trust what the Bayes credential or standard signifies.

The eventual benefit may be strong, but there is an important boundary: \textbf{a private program certificate is not a government-recognized degree or accredited award merely because many colleges use it.} Marketing must preserve that distinction.

\subsection{36. Global ambition starts with standards, not foreign advertisements}

A globally legible institution does not begin by declaring itself global. It begins by creating artifacts that researchers, employers, and technically strong learners anywhere can inspect: rigorous written curricula, challenging open assessments, reproducible projects, technical essays, open-source contributions, faculty biographies, transparent governance, and honest outcomes.

The 1st international learner is useful evidence but does not prove global demand. Overseas acquisition involves different purchasing power, buyer expectations, consumer protection, tax, sanctions, payment processing, data/\allowbreak privacy obligations, accessibility, and time zones. A global cohort can be operationally harder even if teaching occurs on Zoom.

A staged internationalization path:

\begin{itemize}[leftmargin=*,itemsep=1.5pt]
  \item \textbf{Phase G0 - Legibility:} public work meets internationally comprehensible technical standards.
  \item \textbf{Phase G1 - Small demand experiments:} guest instructors, free open seminars, and a few paid international learners from a narrow audience.
  \item \textbf{Phase G2 - Repeatable cross-border program:} proven outcomes, support systems, refund policy, localized payment and compliance review.
  \item \textbf{Phase G3 - Faculty/\allowbreak research network:} coauthorship, visiting appointments, shared research projects, legitimate academic partnerships.
  \item \textbf{Phase G4 - Institutional presence:} exchange programs, permitted formal collaborations, and possibly an overseas center only if justified.
\end{itemize}

A global name is most durable when earned through globally useful work.

-


